\documentclass[a4paper,12pt]{article}

\usepackage[T2A]{fontenc}
\usepackage[utf8]{inputenc}
\usepackage[russian]{babel}
\usepackage{amsmath,amssymb,mathtools}
\usepackage{PTSans}
\renewcommand{\familydefault}{\sfdefault}
\usepackage{icomma}
\usepackage[a4paper,margin=20mm,headheight=25pt]{geometry}
\usepackage{microtype}
\usepackage{tikz}
\usetikzlibrary{arrows.meta,positioning,shapes.geometric,calc}
\usepackage{tabularx,booktabs}
\usepackage{enumitem}
\usepackage[hidelinks]{hyperref}
\usepackage{parskip}
\usepackage{fancyhdr}
\usepackage{setspace}
\usepackage{xcolor}

\definecolor{accent}{HTML}{2A7775}
\definecolor{darkaccent}{HTML}{1B5553}
\definecolor{textgray}{HTML}{303536}
\definecolor{softgray}{HTML}{EEF3F2}

\geometry{a4paper,margin=20mm,headheight=25pt}
\onehalfspacing
\setlength{\parindent}{0pt}
\setlength{\parskip}{0.6em}
\setlist[itemize]{leftmargin=1.35em,itemsep=0.25em,topsep=0.25em}
\setlist[enumerate]{leftmargin=1.7em,itemsep=0.7em,topsep=0.35em}
\hypersetup{colorlinks=true,linkcolor=darkaccent,urlcolor=darkaccent}

\pagestyle{fancy}
\fancyhf{}
\fancyhead[L]{\small\color{textgray}Кирилл Зиновьев}
\fancyhead[R]{\small\color{textgray}07.10.26}
\fancyfoot[C]{\small\color{textgray}\thepage}
\renewcommand{\headrulewidth}{0.3pt}
\renewcommand{\headrule}{\hbox to\headwidth{\color{accent}\leaders\hrule height \headrulewidth\hfill}}

\newcommand{\sectionline}{\par\vspace{0.1em}{\color{accent}\hrule height 0.6pt}\vspace{0.45em}}
\newcommand{\key}[1]{\textbf{\color{darkaccent}#1}}
\newcommand{\R}{\mathbb{R}}

\begin{document}

% ---------------- Title page ----------------
\thispagestyle{empty}
\begin{center}
\vspace*{25mm}

{\Large\color{darkaccent}\textbf{Чек-лист}}\\[8mm]
{\huge\bfseries Уравнения: корни, скобки\\[2mm] и подобные слагаемые}\\[10mm]

{\large 7 класс}\\[3mm]
{\large 07.10.26}\\[15mm]

{\large\textbf{Кирилл Зиновьев}}\\[4mm]
{\normalsize Лёвин Артём Александрович}

\vfill

{\normalsize\color{textgray}Лёвин Артём Александрович эксклюзивно для Кирилла Зиновьева}\\[8mm]

\begin{tikzpicture}
  \draw[accent,line width=1.25pt]
    (0,0) .. controls (2.5,0.9) and (5.0,-0.45) .. (7.5,0.25)
          .. controls (9.2,0.75) and (11.2,0.35) .. (13.0,0.85);
\end{tikzpicture}
\end{center}
\clearpage

% ---------------- Main checklist ----------------
\section*{1. Что значит «корень уравнения»}
\sectionline

\key{Корень уравнения} — это такое значение переменной, при котором левая и правая части уравнения становятся равны.

Чтобы проверить предложенное число, его нужно \key{подставить вместо переменной во все места} и вычислить обе части.

\textbf{Правило оформления:} подставляемое отрицательное число всегда берите в скобки.

\textbf{Пример.} Является ли $x=-3$ корнем уравнения
\[
  x^2+4x+3=0?
\]
Подставляем:
\[
  (-3)^2+4\cdot(-3)+3=9-12+3=0.
\]
Получили ту же величину, что стоит справа, значит $x=-3$ — корень.

\textbf{Самопроверка:}
\begin{itemize}
  \item переменная заменена числом \emph{везде};
  \item отрицательное число записано в скобках;
  \item знаки при умножении посчитаны отдельно;
  \item сравниваются именно левая и правая части.
\end{itemize}

\section*{2. Если произведение равно нулю}
\sectionline

Если
\[
  A\cdot B=0,
\]
то обязательно выполняется хотя бы одно из условий
\[
  A=0 \qquad \text{или} \qquad B=0.
\]

\textbf{Пример.}
\[
  (2x-3{,}8)(4{,}2+3x)=0.
\]
Получаем два простых уравнения:
\[
  2x-3{,}8=0 \quad\Rightarrow\quad x=1{,}9,
\]
\[
  4{,}2+3x=0 \quad\Rightarrow\quad x=-1{,}4.
\]
Следовательно, у исходного уравнения два корня: $1{,}9$ и $-1{,}4$.

\section*{3. Подобные слагаемые}
\sectionline

\key{Подобные слагаемые} имеют одинаковую буквенную часть, включая степени букв.

\begin{center}
\renewcommand{\arraystretch}{1.3}
\begin{tabularx}{0.94\linewidth}{@{}X c X@{}}
\toprule
Пара & Можно объединить? & Почему \\
\midrule
$3x$ и $-5x$ & да & одинаковая буквенная часть $x$ \\
$4a^2$ и $7a^2$ & да & одинаковая буквенная часть $a^2$ \\
$3x$ и $-3$ & нет & у второго слагаемого нет $x$ \\
$3x$ и $-3y$ & нет & разные буквы \\
$3x$ и $-3x^2$ & нет & разные степени \\
\bottomrule
\end{tabularx}
\end{center}

Складывать или вычитать коэффициенты можно только у подобных слагаемых:
\[
  4x+x=5x, \qquad 7a^2-2a^2=5a^2.
\]

\section*{4. Скобки и знак «минус»}
\sectionline

Если перед скобками стоит минус, знак меняется у \key{каждого} слагаемого внутри:
\[
  -(9-x)=-9+x,
\]
\[
  -(a-b-c)=-a+b+c.
\]

\textbf{Полезная привычка:} вычитаемое выражение, подставляемое число и множитель-выражение записывайте в скобках. Это защищает от потери знаков.

\textbf{Пример.}
\[
  5x-1=4(x+2)-(9-x).
\]
Раскрываем скобки:
\[
  5x-1=4x+8-9+x=5x-1.
\]
Левая и правая части совпадают при любом $x$, поэтому корнем является \key{любое действительное число}.

\section*{5. Сколько корней может быть у линейного уравнения}
\sectionline

После раскрытия скобок, переноса и приведения подобных уравнение удобно привести к виду
\[
  ax=b.
\]

\begin{itemize}
  \item Если $a\ne0$, то корень один: $x=\dfrac{b}{a}$.
  \item Если $a=0$ и $b\ne0$, получается невозможное равенство вроде $0=2$ — \key{корней нет}.
  \item Если $a=0$ и $b=0$, получается верное равенство $0=0$ — \key{подходит любое $x\in\R$}.
\end{itemize}

\textbf{Пример без корней.}
\[
  2(x-3)+5x-4=7(x-2)+1.
\]
\[
  7x-10=7x-13 \quad\Rightarrow\quad 0=-3.
\]
Такого быть не может, значит решений нет.

\section*{6. Перевод текста на язык уравнений}
\sectionline

Пусть даны выражения $A$ и $B$.
\[
  \text{сумма: } A+B, \qquad \text{разность: } A-B.
\]

Особенно внимательно записывайте \key{вычитаемое выражение}: оно должно стоять в скобках.

\textbf{Пример 1.} «Сумма выражений $2x+7$ и $-x+12$ равна $24$»:
\[
  (2x+7)+(-x+12)=24.
\]
Отсюда $x+19=24$, поэтому $x=5$.

\textbf{Пример 2.} «Разность выражений $25p+1$ и $p-12$ равна их сумме»:
\[
  (25p+1)-(p-12)=(25p+1)+(p-12).
\]
Главное здесь — поставить скобки вокруг всего вычитаемого выражения.

\clearpage
\section*{7. Короткий контроль перед ответом}
\sectionline

Перед тем как записать ответ, проверьте:
\begin{itemize}
  \item все ли скобки раскрыты корректно;
  \item поменялись ли все знаки после минуса перед скобкой;
  \item объединены ли только подобные слагаемые;
  \item при переносе через знак равенства изменён ли знак;
  \item если переменная сократилась, проверено ли, что осталось: верное или неверное равенство;
  \item ответ соответствует полученному случаю: один корень, нет корней или любое число.
\end{itemize}

\textbf{Три места, где особенно легко ошибиться:}
\[
  x=-3 \quad\Rightarrow\quad (-3)^2,
\]
\[
  -(9-x)=-9+x,
\]
\[
  3x-3x=0, \qquad \text{но } 3x-3 \text{ объединять нельзя}.
\]

Если после упрощения получено $0=0$, подходят все значения переменной. Если получилось $0=c$, где $c\ne0$, решений нет.

\clearpage

% ---------------- Flowchart ----------------
\thispagestyle{fancy}
\begin{center}
{\Large\bfseries\color{darkaccent}Алгоритм: как определить число корней линейного уравнения}
\end{center}
\vspace{8mm}

\begin{center}
\begin{tikzpicture}[
  node distance=11mm and 10mm,
  every node/.style={font=\small,align=center},
  startstop/.style={ellipse,draw=accent,line width=0.9pt,minimum width=40mm,minimum height=11mm},
  process/.style={rectangle,rounded corners=2mm,draw=accent,line width=0.9pt,minimum width=54mm,minimum height=13mm,text width=49mm},
  decision/.style={diamond,aspect=2.1,draw=accent,line width=0.9pt,minimum width=36mm,minimum height=14mm,text width=27mm,inner sep=1.5mm},
  result/.style={rectangle,rounded corners=2mm,draw=darkaccent,line width=1pt,minimum width=38mm,minimum height=13mm,text width=34mm},
  arrow/.style={-{Latex[length=3mm]},line width=0.9pt,color=textgray}
]
  \node[startstop] (start) {Исходное уравнение};
  \node[process,below=of start] (simplify) {Раскрыть скобки, перенести слагаемые, привести подобные};
  \node[process,below=of simplify] (axb) {Получить вид $ax=b$};
  \node[decision,below=of axb] (a0) {$a=0$?};
  \node[result,left=of a0] (one) {$a\ne0$\\[1mm] Один корень\\$x=\dfrac{b}{a}$};
  \node[decision,below=of a0] (b0) {$b=0$?};
  \node[result,left=of b0] (all) {$0=0$\\[1mm] Любое $x\in\R$};
  \node[result,right=of b0] (none) {$0=b$, $b\ne0$\\[1mm] Корней нет};

  \draw[arrow] (start) -- (simplify);
  \draw[arrow] (simplify) -- (axb);
  \draw[arrow] (axb) -- (a0);
  \draw[arrow] (a0.west) -- node[above]{нет} (one.east);
  \draw[arrow] (a0) -- node[right]{да} (b0);
  \draw[arrow] (b0.west) -- node[above]{да} (all.east);
  \draw[arrow] (b0.east) -- node[above]{нет} (none.west);
\end{tikzpicture}
\end{center}

\clearpage

% ---------------- Training ----------------
\section*{Тренировочный блок — Путь А}
\sectionline

Выполните 10 заданий. В задачах на подстановку показывайте саму подстановку; в уравнениях со скобками сначала раскройте скобки, затем приведите подобные.

\begin{enumerate}
  \item Является ли число $x=-2$ корнем уравнения
  \[
    x^2+3x+2=0?
  \]

  \item Какие из чисел $-3$, $-1$, $1$, $3$ являются корнями уравнения
  \[
    x^2+4x+3=0?
  \]

  \item Решите уравнение
  \[
    (3x-6)(2x+5)=0.
  \]

  \item Приведите подобные слагаемые:
  \[
    7x-3+5x+8-2x.
  \]

  \item Выпишите пары подобных слагаемых:
  \[
    4a\text{ и }-9a;\qquad
    3x^2\text{ и }-5x^2;\qquad
    2m\text{ и }2m^2;\qquad
    6p\text{ и }6.
  \]

  \item Решите уравнение
  \[
    4(x-2)+3=4x-5.
  \]

  \item Решите уравнение
  \[
    5(2x+1)-3x=7x+4.
  \]

  \item Решите уравнение
  \[
    3(x-4)+2x=2(x+1)+4.
  \]

  \item Составьте уравнение и найдите $x$: сумма выражений $3x-5$ и $x+1$ равна $20$.

  \item Составьте уравнение и определите множество решений: разность выражений $7p+4$ и $3p-2$ равна сумме выражений $2p+1$ и $2p+5$.
\end{enumerate}

\clearpage

% ---------------- Answers ----------------
\section*{Ответы}
\sectionline

\renewcommand{\arraystretch}{1.45}
\begin{center}
\begin{tabularx}{0.94\linewidth}{@{}c X@{}}
\toprule
\textbf{№} & \textbf{Ответ} \\
\midrule
1 & Да. \\
2 & $-3$ и $-1$. \\
3 & $x=2$; $x=-2{,}5$. \\
4 & $10x+5$. \\
5 & Подобные пары: $4a$ и $-9a$; $3x^2$ и $-5x^2$. \\
6 & Любое $x\in\R$. \\
7 & Корней нет. \\
8 & $x=6$. \\
9 & $x=6$. \\
10 & Любое $p\in\R$. \\
\bottomrule
\end{tabularx}
\end{center}

\vfill
{\small\color{textgray}При самопроверке особенно внимательно смотрите на скобки при подстановке и на смену знаков после минуса перед скобкой.}

\end{document}
